\section{八、产品设计——模型、数据与系统三位一体}

本项目以 BIT 检测模型与专用遥感数据集为核心基础，结合轻量化 Web 前端系统，构建了一套完整的耕地变化监测产品。算法模型解决“看得准”的问题，专用数据集解决“学得会”的问题，可视化前端解决“用得上”的问题，三者协同支撑，具体设计如下。

\subsection{1.核心算法：BIT模型框架}

模型以改进的 ResNet18为 CNN 骨干网络，对双时相输入影像进行特征提取，再通过语义分词器与双时序图像转换器，利用 Transformer 编码器捕捉全局时空上下文信息；最后通过解码器输出精细特征，经预测头生成最终的变化检测结果图，实现了高精度、低误判的黑土地变化检测。

从数据流看，该框架可分为四个阶段：T1、T2两个时刻的影像分别进入共享权重的 ResNet18骨干网络，得到两组空间分辨率逐级降低、通道数逐级增加的中间特征图；语义分词器对特征图做语义级压缩，将其转化为数量可控的语义 token 序列；Transformer 编码器对两个时相的 token 序列进行跨时相交互，借助自注意力在整幅影像范围内建立长程依赖，使同一空间位置在两个时刻的语义特征得以直接比对；解码器再逐级上采样恢复空间分辨率，由预测头输出逐像素的变化概率图。卷积结构与自注意力机制在此形成互补：前者适配遥感影像的局部纹理与边缘特征，后者不受卷积核感受野限制，对形状不规则、分布零散的变化斑块更具判别力，有助于在复杂地表背景下抑制伪变化。项目采用 BIT（Bipartite Transformer）模型\cite{ref1,ref8}作为核心算法，整体框架如下图所示。

四个阶段中，共享权重设计是可比性的前提：两个时相经同一套卷积参数映射到同一特征空间，同一位置的响应才具有统一尺度，否则特征分布出现整体偏移，比对便失去基准。语义分词器承担语义归纳职责，把特征图上数量庞大的空间位置压缩为数量固定的语义 token，每个 token 对应影像中一类语义一致的区域。跨时相交互是判别力的主要来源：编码器把两个时相的 token 序列置于同一注意力计算中，每个 token 都可查询另一时相中语义相近的 token 并交换信息，从而在整幅影像范围内建立长程依赖。解码器再逐级上采样把结论还原到像素级，由预测头输出逐像素的变化概率。

与传统的单时相差分方法相比，双时相交互设计的优势体现在三个方面。其一，差分法直接比较两个时相的像素灰度，对配准误差、光照差异十分敏感，农田翻耕、作物收割等物候现象容易被误判为地表覆盖改变；双时相交互在特征空间中比对语义，学到的是地类是否发生实质改变，对同物异谱、同谱异物现象更为鲁棒。其二，差分法通常只能给出变化与否的二值结论，而跨时相注意力保留了两期影像之间的语义对应关系，可为按类型统计变化面积提供依据。其三，差分阈值需针对不同区域与季节人工标定，迁移到新区域时须重新调参；改进 BIT 的判别边界由数据驱动学习得到，更换监测区域时补充少量样本即可适配。

\begin{figure}[htbp]
    \centering
    \fbox{\parbox[c][4.5cm][c]{0.8\textwidth}{\centering 【图片空位：image9.png】CNN骨干网络逐级提取特征的拓扑结构示意图}}
    \caption{ResNet 骨干网络结构拓扑图}
    \label{fig:5}
\end{figure}

\begin{figure}[htbp]
    \centering
    \fbox{\parbox[c][4.5cm][c]{0.8\textwidth}{\centering 【图片空位：image10.png】BIT双时相变化检测模型整体框架}}
    \caption{BIT模型框架}
    \label{fig:6}
\end{figure}

语义分词器用于从图像特征中提取语义 tokens，配合 Transformer 做图像任务。若把每个空间位置的特征都视为一个 token，注意力计算量将随影像分辨率急剧增长，难以在普通设备上完成推理；语义分词器即在于进入 Transformer 之前完成一次语义压缩，用少量语义明确的 token 概括整幅特征图的信息。

\begin{figure}[htbp]
    \centering
    \fbox{\parbox[c][4.5cm][c]{0.8\textwidth}{\centering 【图片空位：image11.png】语义分词器生成语义token的流程示意}}
    \caption{语义分词器子模块流程图}
    \label{fig:7}
\end{figure}

其中，语义分词器的核心流程如下：

\begin{enumerate}
    \item \textbf{图像特征输入。}输入是维度为 $H \times W \times C$ 的特征图（$H$ 为高度、$W$ 为宽度、$C$ 为通道数），来自 CNN 骨干网络提取的中间特征。
    \item \textbf{卷积降维与特征变换。}用卷积层对特征图进行处理，目的是降维与特征变换，将其映射到适合度量语义相似性的表征空间，为生成注意力图做准备。
    \item \textbf{生成注意力图。}卷积输出经进一步处理得到 $L$ 张注意力图，每张关注特征图中不同的语义区域，用于捕捉关键语义信息，卷积计算过程如图所示。
    \item \textbf{加权聚合得到语义 token。}以注意力图作为空间权重对特征图加权求和，得到 $L$ 个语义 token；其数量与影像分辨率无关，因而在保证语义完整性的同时压缩了后续计算规模。
\end{enumerate}

语义分词器与自注意力机制之间是“先归纳、后关联”的分工：分词器把空间连续、语义一致的区域聚合成少量 token，自注意力则在这些 token 之间建立全局联系。这一分工使注意力计算规模由与像素数相关转为与 token 数相关，而 token 数量可预先设定且不随影像尺寸增长，因此模型面对亚米级高分影像时仍能把显存占用与推理时间控制在可接受范围。注意力图还带来了可解释性：一张注意力图对应影像中一片被模型重点关注的空间区域，多张注意力图叠加后，可以观察模型在判定“发生变化”时关注了哪些位置，为复核人员排查误判提供线索。

\begin{figure}[htbp]
    \centering
    \fbox{\parbox[c][4.5cm][c]{0.8\textwidth}{\centering 【图片空位：image12.png】卷积核在特征图上滑窗计算的过程示意}}
    \caption{卷积计算}
    \label{fig:8}
\end{figure}

该设计使模型在保持识别精度的同时具备良好的推理效率：整体算力成本降低69\%，监测周期缩短33\%，识别精确率达到70\%，为“周级”动态监测提供了算力可行性。

\subsection{2.数据支撑}

算法能力必须建立在高质量数据基础之上。项目构建了黑土地专用遥感数据集，样本以国产高分二号（GF-2）亚米级光学遥感卫星影像为主，辅以无人机航拍影像，覆盖不同季节与地表覆盖类型，以增强模型对真实业务场景的适应能力。表~\ref{tab:dataset}给出该数据集的样本构成。

\begin{table}[htbp]
    \centering
    \caption{黑土地耕地变化检测双时相遥感数据集样本}
    \label{tab:dataset}
    \small
    \begin{tabular}{|c|c|c|c|}
        \hline
        类型 & 时间为T1 & 时间为T2 & 标签 \\
        \hline
        土地变化 & \fbox{\parbox[c][2.2cm][c]{2.6cm}{\centering 【图片空位】\scriptsize image13.png}} & \fbox{\parbox[c][2.2cm][c]{2.6cm}{\centering 【图片空位】\scriptsize image14.jpeg}} & \fbox{\parbox[c][2.2cm][c]{2.6cm}{\centering 【图片空位】\scriptsize image15.png}} \\
        \hline
        海上建设 & \fbox{\parbox[c][2.2cm][c]{2.6cm}{\centering 【图片空位】\scriptsize image16.png}} & \fbox{\parbox[c][2.2cm][c]{2.6cm}{\centering 【图片空位】\scriptsize image17.png}} & \fbox{\parbox[c][2.2cm][c]{2.6cm}{\centering 【图片空位】\scriptsize image18.png}} \\
        \hline
        城市建筑 & \fbox{\parbox[c][2.2cm][c]{2.6cm}{\centering 【图片空位】\scriptsize image19.png}} & \fbox{\parbox[c][2.2cm][c]{2.6cm}{\centering 【图片空位】\scriptsize image20.png}} & \fbox{\parbox[c][2.2cm][c]{2.6cm}{\centering 【图片空位】\scriptsize image21.png}} \\
        \hline
        森林变化 & \fbox{\parbox[c][2.2cm][c]{2.6cm}{\centering 【图片空位】\scriptsize image22.png}} & \fbox{\parbox[c][2.2cm][c]{2.6cm}{\centering 【图片空位】\scriptsize image23.png}} & \fbox{\parbox[c][2.2cm][c]{2.6cm}{\centering 【图片空位】\scriptsize image24.png}} \\
        \hline
    \end{tabular}
\end{table}

数据集构建遵循“采集—配准—标注—核验”的规范化流程：先选取时间基线一致、云覆盖较少、影像质量合格的 T1、T2影像对，完成几何配准与辐射归一化，使两个时相一致对应，排除成像条件差异造成的伪变化；再用 ArcGIS 等专业工具人工标注，逐像素勾绘变化区域边界，形成与影像对一一对应的真值标签；最后通过交叉复核与抽样检查把关质量，对边界模糊、地类易混淆的样本复判或剔除，保证真值可靠性。样本覆盖土地变化、海上建设、城市建筑、森林变化等场景：土地变化样本直接对应耕地被占用、撂荒等地表覆盖改变；海上建设、城市建筑等强变化样本有助于区分真实变化与季节性物候波动；森林变化样本则提升对植被类地物的判别能力。多类型混合训练使模型学习“地表发生变化”的共性规律而非某类地物的外观特征，从而对未见过的变化形态保持泛化能力。

尤需说明的是，数据集的价值不仅在于样本数量，更在于场景分布的代表性。变化检测模型在训练中容易走“捷径”，依据地物外观而非变化本身作出判断，例如把偏绿的影像一律归为植被、把规则矩形一律归为建筑，一旦遇到训练集中未出现的组合便会失效。为此，数据集在构建时有意拉大场景类型的跨度：既包含农田转为建设用地的渐变型变化，也包含地表裸露、水体外扩等突变型变化；既覆盖成片连续的变化，也保留分散、边界不规则的细碎变化。不同类型样本相互约束，迫使模型把判别依据收敛到“两个时刻之间究竟发生了什么”这一核心问题上，以场景覆盖换取对新区域、新变化形态的泛化能力。

需要说明的是，上述场景样本承担的是“通用变化判别”训练任务，其类别设置着眼于覆盖多样的地表变化形态；而黑土地保护业务真正关注的目标类别——侵蚀沟发育、水土流失、土壤有机质流失、耕地违规破坏等——则通过专用样本进行场景适配。二者构成“通用预训练 + 场景适配”的两级训练策略：前者使模型具备对地表变化的通用敏感性，后者把这种敏感性收敛到黑土地退化的具体业务口径上，避免模型把判别依据绑定在某一类地物的外观特征上。

\subsection{3.系统落地：可视化智能前端}

基于前后端分离的轻量化架构，项目开发了一套集影像上传、变化检测、结果可视化、AI 智能解读于一体的现代化遥感检测 Web 前端系统。后端采用 FastAPI 承载模型推理与业务逻辑，前端负责交互呈现，两部分通过标准接口通信，既保证推理服务的稳定性，也便于界面与算法版本的独立迭代。

前后端分离的架构选择同样出于实用考量。算法模型对运行环境存在依赖，若与界面代码耦合，界面调整就可能牵动推理服务重新部署；把推理与呈现拆分为两个独立进程后，模型迭代与前端迭代互不干扰，也便于在不同算力条件的设备上分别优化。

系统支持单张/批量影像上传与处理，自动生成变化检测结果图、热力图与变化比例统计，同时提供一键生成分析报告的功能。单张检测面向定点核查与个案复核，上传一对双时相影像即可获得即时结果；批量检测面向区域性周期性普查，可一次性提交多组影像并统一输出结果与统计汇总，减少人工逐幅操作的工作量。可视化环节给出三类互补信息：结果图以颜色叠加直观标出变化区域的位置与范围；热力图以概率强度反映模型判定的置信程度，帮助区分确定性较高的变化与需人工复核的疑似变化；变化比例统计以结构化方式汇总变化面积占比，为黑土地保护成效评估提供量化依据。AI 智能解读可将像素级输出转译为基层人员易于理解的自然语言描述，一键生成分析报告则将影像、结果图、统计指标与文字解读整合为规范文档，实现从数据输入到决策支撑的全流程闭环。系统面向普通电脑设计，无需专用服务器与高性能显卡即可本地化部署运行，配合算力成本降低69\%、监测周期缩短33\% 的算法优势，降低了一线用户的使用门槛。

前端功能划分遵循“任务—信息—决策”三层递进的取舍原则。任务层只保留影像上传与检测触发两类操作，把参数配置收敛为少量默认项，使不具备遥感背景的基层人员也能独立完成一次检测；信息层在结果图上叠加变化区域、热力图与统计图表，分别回答变化在哪里、模型有多确信、总体变化了多少三个问题；决策层把上述要素汇聚为可下载的分析报告，使检测结果直接进入工作记录与上报流程。单张与批量检测则是两条并列路径：前者面向定点核查，强调上传即得结果，后者面向周期性普查，强调一次提交、统一输出。

综上，模型、数据与系统三者的协同设计，使项目形成了从算法可研到工程可用的完整链路，共同支撑“天空地一体化光学遥感采集 + AI 智能解译 +可视化决策支撑”的业务闭环，服务于黑龙江松嫩平原试点区域的“周级”动态监测需求，并已在松嫩平原试点中服务500+农户、覆盖10万亩黑土地、支撑新增智能监测岗位30+。
